\documentclass[10pt,a4paper]{amsart}
\usepackage[margin=19mm]{geometry}
\usepackage{amsmath,amssymb,mathtools}
\usepackage[scr=rsfs,cal=euler]{mathalfa}
\usepackage{xfrac}
\usepackage[unicode,hidelinks,bookmarks=false]{hyperref}
\newcommand{\ie}{i.e.\ }
\newcommand{\cf}{cf.\ }
\newcommand{\resp}{respectively\ }
\newcommand{\Subsection}{\subsection}
\providecommand{\nobreakdash}{\nobreak\hbox{-}}
\setlength{\emergencystretch}{2em}
\multlinegap0pt
\usepackage{amssymb}
\usepackage{float}
\newtheorem{theorem}{Theorem}
\title{\bf \Large{Green Inventory Management: Leveraging Multiobjective Reverse Logistics}}
\author{I. B. Wadhawan, M. M. Rizvi}
\date{Source: arXiv:2509.19639v2 (CC BY 4.0)}
\begin{document}
\maketitle

\noindent\emph{Source and licence.} \bf \Large{Green Inventory Management: Leveraging Multiobjective Reverse Logistics}. Version arXiv:2509.19639v2 (CC BY 4.0), distributed by arXiv under the Creative Commons Attribution 4.0 International licence, http://creativecommons.org/licenses/by/4.0/, as recorded in the arXiv licence metadata for this exact version and shown on the abstract page https://arxiv.org/abs/2509.19639v2. Full licence notice: \texttt{licenses/arxiv-2509.19639-CC-BY-4.0.txt}.\par\smallskip

\noindent\textit{Editorial scope.} Counted unit: the article's theorem theorem1 (source0.tex lines 350--352). The article's complete original proof of it is delivered with it (source0.tex lines 354--389, one \texttt{proof} environment of 36 lines). No mathematics is written, repaired or re-derived: the statement and the proof are the article's own lines, in the article's own notation, and no retained line is substituted or re-flowed.  Retained uncounted as the source-local context the proof needs: lines 287--289; lines 302--305.  Nothing was omitted from any retained span, so the package carries no omission record.  No retained line contains a citation, so the wrapper carries no bibliography and no bibliography entry.  No retained line contains a figure environment, an image-inclusion command or drawing code, so no image is delivered and the package declares no asset dependency.  Editorial formatting only, in addition: an AMS \texttt{amsart} preamble that restates the article's own declarations and macros used by the retained lines, namely the environments \texttt{theorem} on separate counters, together with the standard packages those lines need.  The retained environments print in this document as Theorem 1 in order of appearance, whereas the article prints the same items with the numbering of its own full document.  The complete native source of the article as distributed in the arXiv e-print for this exact version is bundled unchanged as \texttt{sources/185626/source0.tex}.
\begin{equation}\label{con1}
    mQ_r(1 - sr) + \frac{Q_r}{\lambda} (\lambda - D_r)sr \leq YI = npqQ_p.
\end{equation}

\begin{multline}\label{eqcost1a}
C(Q_p, Q_r, m, n, s) = \frac{1}{T}\Bigg[ mS_n+nS_p + \left(\frac{T_pQ_p(P-D_p)}{2P}+\frac{T_rQ_r(\lambda-D_r)}{2 \lambda}\right)h_p \\ +\Bigg[ \frac{R_1T_p^2}{2} + \frac{m Q_r^2(\lambda - R_2)}{2\lambda^2} +\frac{srQ_rT_r(1-\frac{D_r}{\lambda})^2}{2} \\ 
+\left(1-\frac{1}{m}\right)T_r\left(T_pR_1-\frac{Q_r}{\lambda}(\lambda - R_2)-\frac{(m-2)Q_r(1-sr)}{2}\right)\Bigg] h_r \Bigg],
\end{multline}

\begin{theorem}[Feasibility and Boundedness]\label{theorem1}
Under the above assumptions, the feasible set $X$ is nonempty and bounded.
\end{theorem}

\begin{proof}
\textbf{(Non-emptiness).}
Fix $m=n=1$ and choose any $s\in[0,1]$.
Constraint \eqref{con1} becomes
\[
Q_r(1-sr) + \frac{Q_r}{\lambda}(\lambda-D_r)sr \le pqQ_p.
\]
The left-hand side is linear in $Q_r$, while the right-hand side
is linear in $Q_p$. For any fixed $Q_r\ge1$, choosing $Q_p$
sufficiently large guarantees feasibility. Hence $X\neq\emptyset$.

\medskip

\textbf{(Boundedness).}
Consider the cost function $C(Q_p,Q_r,m,n,s)$ defined in \eqref{eqcost1a}.
The holding cost component contains quadratic growth terms in $Q_p$
and $Q_r$ of the form
\[
\frac{T_pQ_p(P-D_p)}{2P}, \qquad
\frac{T_rQ_r(\lambda-D_r)}{2\lambda},
\]
where $T_p = \frac{nQ_p}{D_p}$ and $T_r = \frac{mQ_r}{D_r}$.

Hence, asymptotically,
\[
C(Q_p,Q_r,m,n,s)
= \Theta(Q_p^2) + \Theta(Q_r^2) + \text{lower-order terms}.
\]
Therefore,
\[
\lim_{\|(Q_p,Q_r,m,n)\|\to\infty} C(Q_p,Q_r,m,n,s) = +\infty.
\]
Since the objective diverges to infinity as decision variables
grow unbounded, any minimizing sequence must lie in a bounded region.
Thus the feasible region containing optimal solutions is bounded.
\end{proof}

\end{document}
